This paper opens a field: Resonant Bioharmonic Medicine (RBM). Rather than treating sound, vibration, and patterned stimulation as peripheral or merely supportive, RBM frames them as candidate physiological interventions with a mechanistic basis and a practical experimental path. The central claim of this work is not that every proposed effect is already established, but that the space of frequency--effect relationships is sufficiently coherent to justify a new research program.

A key strength of RBM is its low barrier to entry. Many of the proposed prototypes and trials can be developed with off-the-shelf hardware, modest software tooling, and careful experimental design. This makes the field accessible to small research teams, independent builders, and translational collaborators who want to test specific hypotheses without waiting for large-scale infrastructure.

Accordingly, this work is an invitation. We are seeking aligned researchers, hardware engineers, and perception-aware investors who are interested in building prototypes, running pilot trials, and refining the measurement frameworks needed to evaluate RBM rigorously. The immediate opportunity is to move from conceptual synthesis to reproducible experimentation, and from isolated observations to a shared technical language for frequency-based intervention.
